\documentclass[12pt]{article}

\usepackage{fancyhdr}
\usepackage{enumitem}
\usepackage{amsmath}
\usepackage{amsfonts}
\usepackage{amssymb}

\renewcommand{\vec}[1]{\mathbf{#1}}

\title{Homework 3}
\author{Satvik Prasad}
\date{}

\pagestyle{fancy}
\lhead{Satvik Prasad}
\rhead{MATH3012}

\begin{document}
\begin{enumerate}
  \item First, let us reduce this problem into the number of integer
    solutions for the equation $$A + B + C' = 13,$$ where $C' = C -
    2.$ The number of solutions to this set of equations can be the
    thought of the number of solutions such that conditions $c_1,$
    $c_2$ and $c_3$ are false, where
    \begin{align*}
      c_1 : A \geq 7, \quad c_2 : B \geq 12, \quad c_3 : C' \geq 6.
    \end{align*}
    Using inclusion exclusion, we know
    \begin{align*}
      N(\bar c_1 \bar c_2 \bar c_3) =  S_0 - S_1 + S_2 - S_3.
    \end{align*}
    In this case, $S_0$ is the number of unique non-negative integer
    solutions, which can be found by distributing $13$ stars over $2$
    bars, meaning $$S_0 = {15 \choose 2}.$$
    Similarly, $S_1$ represents the sum of all the ways to have one
    of the conditions present, which can be enumerated:
    \begin{itemize}
      \item If $c_1,$ then we can fix $7$ stars into bucket $A$, and
        distribute $6$ stars over $2$ bars, yielding ${8 \choose 2}$ ways.
      \item If $c_2,$ in a similar fashion we fix $12$ stars into
        bucket $B,$ leaving $1$ star to distribute over $2$ bars,
        yielding ${3 \choose 1}$ ways.
      \item If $c_3,$ then we can fix $6$ stars into bucket $C',$
        distributing $7$ stars over $2$ bars, yielding ${9 \choose 2}$ ways.
    \end{itemize}
    Hence, $S_1 = {8 \choose 2} + {3 \choose 1} + {9 \choose 2}.$

    Thirdly, $S_2$ represents the sum of ways to have pairs of
    conditions true. If $c_1$ and $c_2,$ then $A + B \geq 19$
    yielding $0$ possibilities. If $c_1$ and $c_3,$ $A + C' = 13,$
    meaning $B$ must be 0, yielding $1$ possibility. Finally, if
    $c_2$ and $c_3$, then $B + C' \geq 18,$ yielding $0$
    possibilities. Hence, $S_2 = 1.$

    Finally, $S_3$ represents the ways to have all three conditions
    true. In this case, $A \geq 7, B \geq 12, C' \geq 6$ meaning
    their sum is greater than or equal to 25, yielding $0$ possible
    solutions. Hence,
    \begin{align*}
      N(\bar c_1 \bar c_2 \bar c_3) &= {15 \choose 2} - \left({8
      \choose 2} + {3 \choose 1} + {9 \choose 2}\right) + 1 - 0 \\
      &= 105 - (28 + 3 + 36) + 1 \\
      &= 39
    \end{align*}
    possible unique integer solutions.

    \iffalse
    ARRANGEMENT

    ARNGEMT
    ARN E
    \fi

  \item
    \begin{enumerate}
      \item The singletons in ARRANGEMENT are G, M, and T, whereas
        the letters that occur twice are A, R, N, and E.

        Let $E_2$ represent the number of ways to have exactly two
        pairs of consecutive identical letters. Then, using
        inclusion-exclusion, we know
        \begin{align*}
          E_2 = S_2 - {3 \choose 2}S_3 + {4 \choose 2}S_4.
        \end{align*}

        We can compute $S_2$ by choosing 2 out of the 4 letters to be
        adjacent, treating adjacent pairs as a single unit, and
        permuting, yielding
        \begin{align*}
          S_2 = {4 \choose 2} \cdot \dfrac{9!}{(2!)^2}.
        \end{align*}
        Similarly, we can compute $S_3$ by choosing $3$ out of the
        $4$ letters to come in adjacent pairs, and then permuting:
        \begin{align*}
          S_3 = {4 \choose 3} \cdot \dfrac{8!}{2!}.
        \end{align*}
        Finally, we can compute $S_4$ by choosing all the letters to
        come in adjacent pairs:
        \begin{align*}
          S_4 = 7!.
        \end{align*}
        Hence,
        \begin{align*}
          E_2 &= {4 \choose 2} \dfrac{9!}{(2!)^2} - {3 \choose 2}
          \cdot {4 \choose 3} \cdot \dfrac{8!}{2!} + {4 \choose 2} \cdot 7! \\
          &= 544320 - 241920 + 30240 \\
          &= 332640
        \end{align*}
        ways.

      \item We can use the fact that
        \begin{align*}
          L_2 &= S_0 - E_1 - E_0 \\
          &= S_0 - \left(S_1 - {2 \choose 1}S_2 + {3 \choose 1}S_3 -
          {4 \choose 1}S_4\right) \\ &- \left(S_0 - S_1 + S_2 - S_3 +
          S_4\right) \\
          &= S_2 - 2S_3 + 3S_4.
        \end{align*}
        Using our values for $S_k$ in the previous question, we get
        \begin{align*}
          L_2 &= {4 \choose 2} \cdot \dfrac{9!}{(2!)^2   } - 2\cdot
          {4 \choose 3} \cdot \dfrac{8!}{2!} + 3 \cdot 7! \\
          &= 544320 - 161280 + 15120 \\
          &= 398160
        \end{align*}
        total ways to have at least two pairs of consecutive identical letters.
    \end{enumerate}

    \iffalse
    Determine derangement of n elements. All elements must be in
    different spots. Consider first element, which has $n - 1$
    locations to go. So $n - 1$ ways to place it. Then, if the next
    number goes in the first spot, we derange the rest $D(n - 2)$
    ways. Otherwise, we treat the start as the place where the chosen
    number cannot go, yielding $D(n - 1)$ ways.
    \fi

  \item
    \begin{enumerate}
      \item We first determine an arrangement for the first
        classroom, and then a derangement for the second classroom.
        To arrange the first classroom, we have $12!$ ways. Then, to
        derange the 12 students in the second classroom, we have
        $$D(12) = 12! \sum_{k = 0}^{12}\dfrac{(-1)^k}{k!} =
        \left[\dfrac{12!}{e}\right] = 176214841.$$ Hence, we have a total of
        \begin{align*}
          (12!)^2 \sum_{k = 0}^{12}\dfrac{(-1)^k}{k!} = 12! \cdot (176214841)
        \end{align*} ways.

      \item We choose the $6$ students to occupy the same desks,
        arrange the first classroom, and then derange the other $6$
        students in the second classroom, yielding
        \begin{align*}
          &{12 \choose 6} \cdot 12! \cdot D(6) \\
          &= \left(\dfrac{12!}{6!}\right)^2 6! \sum_{k = 0}^6
          \dfrac{(-1)^k}{k!}. \\
          &= \dfrac{(12!)^2}{6!} \sum_{k = 0}^{6} \dfrac{(-1)^k}{k!} \\
          &= \dfrac{(12!)^2}{(6!)^2} \left[\dfrac{6!}{e}\right] \\
          &= \dfrac{(12!)^2}{(6!)^2} \cdot 265
        \end{align*}
        ways.
    \end{enumerate}

\end{enumerate}
\end{document}
